% Generated draft: requires mathematical and source review.
\documentclass[11pt]{article}
\usepackage[margin=25mm]{geometry}
\usepackage{fontspec}
\setmainfont{Latin Modern Roman}
\usepackage{amsmath,amssymb,mathtools,unicode-math}
\setmathfont{Latin Modern Math}
\usepackage{enumitem,needspace,xurl,hyperref}
\setlength{\parindent}{0pt}
\setlength{\parskip}{5pt}
\setlength{\emergencystretch}{3em}
\newcommand{\N}{\mathbb{N}}
\newcommand{\Z}{\mathbb{Z}}
\newcommand{\Q}{\mathbb{Q}}
\newcommand{\R}{\mathbb{R}}
\newcommand{\C}{\mathbb{C}}
\newcommand{\F}{\mathbb{F}}
\newcommand{\A}{\mathbb{A}}
\newcommand{\PP}{\mathbb{P}}
\newcommand{\E}{\mathbb{E}}
\newcommand{\eps}{\varepsilon}
\newcommand{\dd}{\,\mathrm{d}}
\begin{document}
\section*{UnsolvedMath ID: 30003165}
\textbf{Source problem code:} catalogue entry
\subsection{Definitions and mathematical statement}
Write \(M(T,B)\) for the motive with modulus attached to a modulus pair \((T,B)\) in the intended source formalism, and write \(|B|\) for the support of its divisor. Under the usual effective Cartier divisor convention, such a pair consists of an ambient scheme \(T\) and an effective Cartier divisor \(B\) with smooth complement \(T\setminus |B|\); the base, finiteness conditions, and precise category of modulus pairs must still be specified. This convention is explanatory and does not resolve the missing source hypotheses. Write \(CH^d(V,j)\) for the ordinary higher Chow group of codimension \(d\) and degree \(j\), and \([-j]\) for translation in the intended triangulated category.

The supplied starting formula is
\[
\operatorname{Hom}\bigl(M(Y,E),M(X,0)[-j]\bigr)
\cong CH^d\bigl(X\times(Y\setminus |E|),j\bigr),
\]
where \(X\) is smooth, proper, and equidimensional of dimension \(d\), and \((Y,E)\) is a modulus pair. The product is over the common base, which the record does not specify. Neither the category in which \(\operatorname{Hom}\) is taken, the coefficient convention, nor the range of \(j\) is supplied; the displayed formula is retained as the source's starting assertion with these qualifications.

\begin{enumerate}
\item For such \(X\) and \((Y,E)\), and a nonzero divisor \(F\) on \(X\) admissible as a modulus in the intended formalism, formulate and establish an analogue of the displayed isomorphism with target \(M(X,F)[-j]\). The source suggests that higher Chow groups with modulus may provide the right-hand side. The requested formula should recover the displayed formula when the target modulus is zero. No particular higher Chow group with modulus, modulus condition on a product, or rule relating \(E\) and \(F\) is specified, so a definite right-hand side cannot be written from this record alone.
\item For a curve \(C\) with modulus \(D\), describe explicitly a complex \(K^\bullet\) of objects of \(\mathrm{MNST}\) representing the indicated object, that is, give an isomorphism
\[
K^\bullet\simeq i\bigl(M(C,D)\bigr)
\quad\text{in }D(\mathrm{MNST}).
\]
Here \(\mathrm{MNST}\), \(D(\mathrm{MNST})\), and \(i\) retain their source meanings. Their precise definitions, the domain and construction of \(i\), and any boundedness convention are not supplied. No additional smoothness, properness, or connectedness assumptions on \(C\) are stated in the record.
\end{enumerate}

These are formulation and explicit-description questions. The missing categorical and geometric conventions remain unresolved; the record does not determine a fully specified theorem.
\subsection{Short English statement}
The problem asks for an extension of a supplied motivic Hom--higher Chow group formula to a nonzero target modulus, possibly using higher Chow groups with modulus. It also asks for an explicit complex representing the motive of a curve with modulus after applying \(i\). The base, categories, coefficients, degree range, and functor \(i\) require clarification before either question becomes fully precise.
\subsection{Sources}
\begin{itemize}
\item Supplied local UnsolvedMath registry record, permanent ID \(30003165\), title \emph{Motives and Higher Chow Groups with Modulus}, source problem code \texttt{catalogue entry}, in the Oberwolfach Reports Open Problems collection. The record gives proposed year 2016 and supplies no individual proposer.
\item Supplied report citation: \emph{Algebraic K-theory and Motivic Cohomology} (2017), Oberwolfach Reports, pp.\ 1753--1807, DOI 10.4171/owr/2016/31. \url{https://doi.org/10.4171/owr/2016/31}. The collection describes these reports as records of mathematical workshops and their open-problem discussions. The report contents have not been independently inspected for this draft.
\item Supplied literature source: Kahn--Saito--Yamazaki, \emph{Reciprocity sheaves}. \url{https://arxiv.org/abs/1402.4201}.
\item Supplied literature source: \emph{Reciprocity sheaves and logarithmic motives}. \url{https://arxiv.org/abs/2107.00381}.
\item Literature assessment date: \texttt{2026-08-21}. This is supplied dated curation, not an independently verified assessment. The record labels the problem open and reports that later reciprocity-sheaf and logarithmic-motive work develops the surrounding framework, but that no primary result located through that date gives the requested nontrivial-divisor formula and explicit complex in the stated derived category. The supplied curation calls for independent verification of time-sensitive status and cited claims.
\end{itemize}
\end{document}
